\subsection{Class \texttt{trajectory}}
\label{sec:trajectory}

Trajectory data is stored in CORA as objects of the class \texttt{trajectory}, which provides several methods to easily visualize the simulated trajectories. An object of class \texttt{trajectory} can be constructed as follows:
\begin{equation*}
	\begin{split}
		& \texttt{traj} = \texttt{trajectory}(\texttt{u},\texttt{x},\texttt{y},\texttt{t}), \\
		& \texttt{traj} = \texttt{trajectory}(\texttt{u},\texttt{x},\texttt{y},\texttt{[]},\texttt{dt}), \\
		& \texttt{traj} = \texttt{trajectory}(\texttt{u},\texttt{x},\texttt{y},\texttt{t},\texttt{dt},\texttt{loc}), \\
		& \texttt{traj} = \texttt{trajectory}(\texttt{u},\texttt{x},\texttt{y},\texttt{t},\texttt{dt},\texttt{[]},\texttt{a}), \\
		& \texttt{traj} = \texttt{trajectory}(\texttt{u},\texttt{x},\texttt{y},\texttt{t},\texttt{dt},\texttt{loc},\texttt{a},\texttt{name}), \\
		& \texttt{traj} = \texttt{trajectory}(\texttt{u},\texttt{x},\texttt{y},\texttt{t},\texttt{dt},\texttt{loc},\texttt{a},\texttt{name},\texttt{model}),
	\end{split}
\end{equation*}
with input arguments
\begin{center}
\renewcommand{\arraystretch}{1.3}
\begin{tabular}[t]{l p{13cm} }
	$\bullet$~\texttt{u} & array storing the input of the trajectory, where the number of rows is identical to the number of
	system inputs and the number of columns is identical to the length of the trajectory.\\
	$\bullet$~\texttt{x} & array storing the states of the trajectory, where the number of rows is identical to the number of
	system states and the number of columns is identical to the length of the trajectory.\\
	$\bullet$~\texttt{y} & array storing the outputs of the trajectory, where the number of rows is identical to the number of
	system outputs and the number of columns is identical to the length of the trajectory.\\
	$\bullet$~\texttt{t} & row vector storing the time points for the trajectory. \\
	$\bullet$~\texttt{dt} & time step size.\\
	$\bullet$~\texttt{loc} & array storing the indices of the locations for the trajectory (nonempty for hybrid systems only), where the number of rows is identical to the number of
	locations and the number of columns is identical to the length of the trajectory.\\
	$\bullet$~\texttt{a} & array storing algebraic variables for the trajectory (for differential-algebraic systems), where the number of rows is identical to the number of
	algebraic variables and the number of columns is identical to the length of the trajectory. \\
	$\bullet$~\texttt{name} & name for the trajectory. \\
	$\bullet$~\texttt{model} & function handle for the model. \\
\end{tabular}
\end{center}

Next, we explain some methods of the class \texttt{trajectory} in detail.

\subsubsection{\texttt{add}}

The method \texttt{add} combines two trajectories \texttt{traj1} and \texttt{traj2} by concatenating the states, outputs, and algebraic variables along the third dimension if \texttt{traj1} and \texttt{traj2} have the same dimensions, inputs, initial state, time vector, and locations.
\begin{equation*}
	\texttt{traj} = \texttt{add}(\texttt{traj1},\texttt{traj2}).
\end{equation*}


\subsubsection{\texttt{plot}}

The method \texttt{plot} visualizes a two-dimensional projection of the obtained trajectories:

\begin{equation*}
	\begin{split}
		&\texttt{han} = \operator{plot}(\texttt{traj}), \\
		&\texttt{han} = \operator{plot}(\texttt{traj},\texttt{dims}), \\
		&\texttt{han} = \operator{plot}(\texttt{traj},\texttt{dims},\texttt{linespec}), \\
		&\texttt{han} = \operator{plot}(\texttt{traj},\texttt{dims},\texttt{namevaluepairs}),
	\end{split}
\end{equation*}
where \texttt{traj} is an object of class \texttt{trajectory}, \texttt{han} is a handle to the plotted MATLAB graphics object, and the additional input arguments are defined as

\begin{itemize}
	\item \texttt{dims}: Integer vector $\texttt{dims} \in \mathbb{N}_{\leq n}^2$ specifying the dimensions for which the projection is visualized (default value: \texttt{dims = [1 2]}).
	\item \texttt{linespec}: (optional) line specifications, e.g., \texttt{'--*r'}, as supported by MATLAB\footnote{\url{https://de.mathworks.com/help/matlab/ref/linespec.html}} (default value: \texttt{linespec = 'b'}).
	\item \texttt{namevaluepairs}: (optional) further specifications as name-value pairs, e.g., \texttt{'Traj','y'} to specify which trajectory to plot (default is \texttt{'Traj','x'}), or line properties such as
	\texttt{'LineWidth',2} and \texttt{'MarkerSize',1.5}, as supported by MATLAB.
	They correspond to the Line Properties\footnote{\url{https://de.mathworks.com/help/matlab/ref/matlab.graphics.chart.primitive.line-properties.html}}.
\end{itemize}


\subsubsection{\texttt{plotOverTime}}

The method \texttt{plotOverTime} visualizes a one-dimensional projection of the simulated trajectories over time:

\begin{equation*}
	\begin{split}
		&\texttt{han} = \operator{plotOverTime}(\texttt{traj}), \\
		&\texttt{han} = \operator{plotOverTime}(\texttt{traj},\texttt{dims}), \\
		&\texttt{han} = \operator{plotOverTime}(\texttt{traj},\texttt{dims},\texttt{linespec}), \\
		&\texttt{han} = \operator{plotOverTime}(\texttt{traj},\texttt{dims},\texttt{namevaluepairs}),
	\end{split}
\end{equation*}
where \texttt{traj} is an object of class \texttt{trajectory}, \texttt{han} is a handle to the plotted MATLAB graphics object, and the additional input arguments are defined as

\begin{itemize}
	\item \texttt{dims}: Integer vector $\texttt{dims} \in \mathbb{N}_{\leq n}$ specifying the dimensions for which the projection is visualized (default value: \texttt{dims = 1}).
	\item \texttt{linespec}: (optional) line specifications, e.g., \texttt{'--*r'}, as supported by MATLAB\footnote{\url{https://de.mathworks.com/help/matlab/ref/linespec.html}}.
	\item \texttt{namevaluepairs}: (optional) further specifications as name-value pairs, e.g., \texttt{'Traj','y'} to specify which trajectory to plot (default is \texttt{'Traj','x'}), or line properties such as
	\texttt{'LineWidth',2} and \texttt{'MarkerSize',1.5}, as supported by MATLAB.
	They correspond to the Line Properties\footnote{\url{https://de.mathworks.com/help/matlab/ref/matlab.graphics.primitive.patch-properties.html}}.
\end{itemize}
